\section{去偏估计器的证明}\label{app:estimator}

本附录固定一个冻结目标，简记$g=\nabla J(\theta)$。在相应条件分布下，完整轨迹摘要$\mathsf X_1,\mathsf X_2,\ldots$独立同分布。对样本$x=(V^+,V^-,G_x)$，记$I_{x,\sigma}(y)=\1\{V^\sigma>y\}$。

\subsection{二阶展开与全样本/半样本耦合}

\begin{lemma}[耦合的全样本/半样本抵消]\label{lem:level-cancellation}
在假设~\ref{ass:trajectory}--\ref{ass:cpt-regularity}下，固定冻结目标与参数，令$g:=\nabla J(\theta)$。存在可测函数$\mathfrak h$及有限常数$C_{\rm rem},C_U$，使
\begin{equation}\label{eq:loo-expansion}
 U_n^{\rm raw}-g
 =\frac1n\sum_{i=1}^{n}\mathfrak h(\mathsf X_i)+\mathcal B_n,
 \qquad
 \E\mathfrak h(\mathsf X_i)=0,
 \qquad
 \E\norm{\mathcal B_n}_2^2\le\frac{C_{\rm rem}}{n^2},
\end{equation}
从而
\begin{equation}\label{eq:level-second-moment}
 \E\norm{D_{n,\ell}^{\rm raw}}_2^2
 \le\frac{9C_{\rm rem}}{n^2}2^{-2\ell},
 \qquad \ell\ge1,
\end{equation}
且
\begin{equation}\label{eq:centered-loo-variance}
 \tr\Cov(U_n^{\rm cen})\le\frac{C_U}{n}.
\end{equation}
\end{lemma}

\begin{proof}
先从$U_n^{\rm raw}$中分离一阶Hoeffding投影，并分别控制退化双样本项与Taylor余项。对侧别$\sigma$和阈值$y$，记
\[
 S_\sigma(y)=\Pp(V^\sigma>y),
 \qquad
 e_{-i,\sigma}(y)=\widehat S_{-i,n}^\sigma(y)-S_\sigma(y).
\]
对$(w_\sigma^\eps)'$应用Taylor定理，得到
\begin{equation}\label{eq:weight-taylor}
 (w_\sigma^\eps)'(S_\sigma+e)
 =(w_\sigma^\eps)'(S_\sigma)
 +(w_\sigma^\eps)''(S_\sigma)e
 +\operatorname{Rem}_\sigma(e),
 \qquad
 \abs{\operatorname{Rem}_\sigma(e)}\le\frac{C_3}{2}e^2.
\end{equation}
把式~\eqref{eq:weight-taylor}代入式~\eqref{eq:raw-loo}，统计量分为一阶轨迹项、按顺序计数的双样本项和Taylor余项。

定义
\begin{align*}
 \mathfrak h_0(x)
 &=G_x\sum_\sigma \varsigma_\sigma
 \int_0^{B_v}(w_\sigma^\eps)'(S_\sigma(y))I_{x,\sigma}(y)\,dy,\\
 \mathcal K(x,x')
 &=G_x\sum_\sigma \varsigma_\sigma
 \int_0^{B_v}(w_\sigma^\eps)''(S_\sigma(y))
 I_{x,\sigma}(y)\{I_{x',\sigma}(y)-S_\sigma(y)\}\,dy.
\end{align*}
由式~\eqref{eq:score-gradient-identity}和$\E G=0$，有$\E\mathfrak h_0(\mathsf X)=g$。对独立的$\mathsf X,\mathsf X'$，
$\E\{\mathcal K(\mathsf X,\mathsf X')\mid\mathsf X\}=0$，因为$\E\{I_{\mathsf X',\sigma}(y)-S_\sigma(y)\}=0$。令
\[
 \mathfrak h_1(x')=\E_{\mathsf X}\{\mathcal K(\mathsf X,x')\}.
\]
则$\E\mathfrak h_1(\mathsf X')=0$。定义$\mathcal K^\circ(x,x')=\mathcal K(x,x')-\mathfrak h_1(x')$后，其两个一阶投影均为零：
\[
 \E\{\mathcal K^\circ(\mathsf X,\mathsf X')\mid\mathsf X\}=0,
 \qquad
 \E\{\mathcal K^\circ(\mathsf X,\mathsf X')\mid\mathsf X'\}=0.
\]
所以双样本项具有精确分解
\begin{equation}\label{eq:ordered-hoeffding}
 \frac{1}{n(n-1)}\sum_{i\ne j}\mathcal K(\mathsf X_i,\mathsf X_j)
 =\frac1n\sum_{j=1}^{n}\mathfrak h_1(\mathsf X_j)
 +\mathcal D_n,
\end{equation}
其中
\[
 \mathcal D_n=\frac{1}{n(n-1)}\sum_{i\ne j}\mathcal K^\circ(\mathsf X_i,\mathsf X_j).
\]
指标函数有界、积分区间有限、$(w_\sigma^\eps)''$有界，且$\E\norm G_2^2<\infty$，故核函数平方可积。退化性使两项的协方差只有在其有序下标对相同或相反时才可能非零。这样的下标对只有$O(n^2)$个，而分母平方阶为$n^2(n-1)^2$，因此存在有限$C_D$，使
\begin{equation}\label{eq:degenerate-remainder}
 \E\norm{\mathcal D_n}_2^2\le\frac{C_D}{n^2}.
\end{equation}

再考察Taylor余项。对每个$i$，$e_{-i,\sigma}$是由$n-1$个与$\mathsf X_i$独立的伯努利变量构成的经验均值。对所有$y$，有
\begin{equation}\label{eq:bernoulli-fourth-moment}
 \E\abs{e_{-i,\sigma}(y)}^4\le\frac{C_e}{(n-1)^2},
\end{equation}
其中$C_e$是通用有限常数。由式~\eqref{eq:weight-taylor}、$y$积分上的Cauchy--Schwarz不等式、$G_i$与留一经验生存概率的独立性及得分二阶矩有限，每个Taylor项的二阶矩至多为$C_T/n^2$。再对$i$的平均使用Jensen不等式，得到平均余项$\mathcal T_n$的相同阶界：
\begin{equation}\label{eq:taylor-remainder-bound}
 \E\norm{\mathcal T_n}_2^2\le\frac{C_T}{n^2}.
\end{equation}

合并一阶项、式~\eqref{eq:ordered-hoeffding}和Taylor余项，
\[
 U_n^{\rm raw}-g
 =\frac1n\sum_{i=1}^{n}\mathfrak h(\mathsf X_i)+\mathcal B_n,
 \qquad
 \mathfrak h=\mathfrak h_0-g+\mathfrak h_1,
 \qquad
 \mathcal B_n=\mathcal D_n+\mathcal T_n.
\]
$\mathfrak h$的期望为零；放大有限常数$C_{\rm rem}$后，由式~\eqref{eq:degenerate-remainder}--\eqref{eq:taylor-remainder-bound}得到$\E\norm{\mathcal B_n}_2^2\le C_{\rm rem}/n^2$，从而证明式~\eqref{eq:loo-expansion}。

现在取$n_\ell=n2^\ell$，在同一批$n_\ell$条轨迹上计算全样本及两个半样本的未中心化统计量。对全样本与各半样本应用式~\eqref{eq:loo-expansion}，有
\[
 U_{n_\ell}^{\rm raw}
 =g+\frac1{n_\ell}\sum_{i=1}^{n_\ell}\mathfrak h(\mathsf X_i)+\mathcal B_{n_\ell},
\]
以及
\[
 \frac12\{U_{n_\ell/2}^{\rm raw,(1)}+U_{n_\ell/2}^{\rm raw,(2)}\}
 =g+\frac1{n_\ell}\sum_{i=1}^{n_\ell}\mathfrak h(\mathsf X_i)
 +\frac12\{\mathcal B_{n_\ell/2}^{(1)}+\mathcal B_{n_\ell/2}^{(2)}\}.
\]
一阶影响项完全抵消，留下
\[
 D_{n,\ell}^{\rm raw}
 =\mathcal B_{n_\ell}-\frac12\{\mathcal B_{n_\ell/2}^{(1)}+\mathcal B_{n_\ell/2}^{(2)}\}.
\]
由$L^2$中的Minkowski不等式，
\begin{align*}
 (\E\norm{D_{n,\ell}^{\rm raw}}_2^2)^{1/2}
 &\le\frac{\sqrt{C_{\rm rem}}}{n_\ell}
 +\frac12\frac{2\sqrt{C_{\rm rem}}}{n_\ell}
 +\frac12\frac{2\sqrt{C_{\rm rem}}}{n_\ell}
 =\frac{3\sqrt{C_{\rm rem}}}{n_\ell}.
\end{align*}
两边平方并代入$n_\ell=n2^\ell$，得到式~\eqref{eq:level-second-moment}。

最后证明式~\eqref{eq:centered-loo-variance}。将$U_n^{\rm cen}$视为$\mathsf X_1,\ldots,\mathsf X_n$的对称函数，用独立副本$\mathsf X_i'$替换第$i$个样本后的统计量记为$U_n^{\rm cen,(i)}$。若$G_i'$为该副本的得分，则单个中心化留一项的系数绝对值不超过$2B_vC_1$，被替换项的贡献范数至多为
\[
 \frac{2B_vC_1}{n}(\norm{G_i}_2+\norm{G_i'}_2).
\]
对$j\ne i$，替换$\mathsf X_i$使第$j$项的留一生存概率逐点变化不超过$(n-1)^{-1}$。由于
\[
 \left|\frac{\partial}{\partial p}\{(w_\sigma^\eps)'(p)(I-p)\}\right|
 \le C_1+C_2,
\]
乘在$G_j$前的标量系数变化不超过$2B_v(C_1+C_2)/(n-1)$。因而存在只依赖$B_v,C_1,C_2$的确定常数$C_*$，使
\begin{align*}
 \norm{U_n^{\rm cen}-U_n^{\rm cen,(i)}}_2
 &\le \frac{C_*}{n}(\norm{G_i}_2+\norm{G_i'}_2)
 +\frac{C_*}{n(n-1)}\sum_{j\ne i}\norm{G_j}_2.
\end{align*}
使用$(\sum_{j\ne i}x_j)^2\le(n-1)\sum_{j\ne i}x_j^2$和得分的统一二阶矩界，可取有限的$C_U$，使对所有$i$都有
\[
 \E\norm{U_n^{\rm cen}-U_n^{\rm cen,(i)}}_2^2
 \le\frac{2C_U}{n^2}.
\]
向量形式的Efron--Stein不等式遂给出
\[
 \tr\Cov(U_n^{\rm cen})
 =\E\norm{U_n^{\rm cen}-\E U_n^{\rm cen}}_2^2
 \le\frac12\sum_{i=1}^{n}
 \E\norm{U_n^{\rm cen}-U_n^{\rm cen,(i)}}_2^2
 \le\frac{C_U}{n}.
\]
证毕。
\end{proof}

\subsection{中心化与有限样本偏差}

\begin{lemma}[中心化与有限样本偏差]\label{lem:centering-expectation}
对$n\ge2$，有
\begin{equation}\label{eq:centered-raw-mean}
 \E U_n^{\rm cen}=\E U_n^{\rm raw}=: \bar U_n,
\end{equation}
且
\begin{equation}\label{eq:loo-bias-order}
 \norm{\bar U_n-g}_2\le \frac{\sqrt{C_{\rm rem}}}{n},
 \qquad g:=\nabla J(\theta).
\end{equation}
\end{lemma}

\begin{proof}
由式~\eqref{eq:raw-loo}减去式~\eqref{eq:centered-loo}，得到
\begin{equation}\label{eq:centering-control}
 U_n^{\rm raw}-U_n^{\rm cen}
 =\frac1n\sum_{i=1}^{n}G_i H_{-i}^{\rm ctr},
\end{equation}
其中
\[
 H_{-i}^{\rm ctr}
 =\sum_{\sigma\in\{+,-\}}\varsigma_\sigma
 \int_0^{B_v}
 (w_\sigma^\eps)'(\widehat S_{-i,n}^\sigma(y))
 \widehat S_{-i,n}^\sigma(y)\,dy.
\]
$H_{-i}^{\rm ctr}$只依赖下标不为$i$的轨迹，因而与$G_i$独立。得分恒等式给出$\E G_i=0$，而假设~\ref{ass:trajectory}--\ref{ass:cpt-regularity}保证$G_iH_{-i}^{\rm ctr}$可积。所以式~\eqref{eq:centering-control}中的每一项期望均为零，得到$\E U_n^{\rm cen}=\E U_n^{\rm raw}=\bar U_n$。

再由引理~\ref{lem:level-cancellation}中的展开式~\eqref{eq:loo-expansion}及$\E\mathfrak h(\mathsf X)=0$，
\[
 \norm{\bar U_n-g}_2
 =\norm{\E \mathcal B_n}_2
 \le(\E\norm{\mathcal B_n}_2^2)^{1/2}
 \le\frac{\sqrt{C_{\rm rem}}}{n}.
\]
这即是式~\eqref{eq:loo-bias-order}。
\end{proof}

\subsection{无偏性、方差、成本与中心极限定理}

\begin{proof}
先证明严格无偏。由引理~\ref{lem:centering-expectation}，中心化基项的期望为$\bar U_n$。未中心化修正项满足
\begin{equation}\label{eq:delta-mean}
 \E D_{n,\ell}^{\rm raw}
 =\bar U_{n2^\ell}-\bar U_{n2^{\ell-1}}.
\end{equation}
式~\eqref{eq:level-second-moment}推出
\[
 \sum_{\ell\ge1}\E\norm{D_{n,\ell}^{\rm raw}}_2
 \le\frac{3\sqrt{C_{\rm rem}}}{n}\sum_{\ell\ge1}2^{-\ell}<\infty,
\]
故可交换期望与无限随机级数。又由式~\eqref{eq:loo-bias-order}，$\bar U_{n2^{\ell_{\max}}}\to g$，于是
\begin{align*}
 \E Z_n^{\db}
 &=\bar U_n+\sum_{\ell\ge1}\E D_{n,\ell}^{\rm raw}\\
 &=\bar U_n+\lim_{\ell_{\max}\to\infty}
 (\bar U_{n2^{\ell_{\max}}}-\bar U_n)
 =g.
\end{align*}
这证明式~\eqref{eq:estimator-unbiased}。

激活变量、层级变量与修正样本独立，因此修正项的二阶矩为
\begin{align*}
 \E\left\|\frac{\mathsf A}{p_{\rm act}\nu_{\ell^\star}}D_{n,\ell^\star}^{\rm raw}\right\|_2^2
 &=\frac{1}{p_{\rm act}}\sum_{\ell\ge1}
 \frac{\E\norm{D_{n,\ell}^{\rm raw}}_2^2}{\nu_\ell}\\
 &\le\frac{9C_{\rm rem}}{p_{\rm act} n^2}
 \sum_{\ell\ge1}\frac{4^{-\ell}}{\nu_\ell}
 =\frac{9C_{\rm rem}\Xi_b}{p_{\rm act} n^2}.
\end{align*}
因$b<2$，级数$\Xi_b$有限。基项与修正项独立，协方差为零；结合式~\eqref{eq:centered-loo-variance}，得到式~\eqref{eq:single-output-variance}。

一次基项输出恰好使用$n$条完整轨迹。若修正被激活且抽到层级$\ell$，只需额外模拟一批$n2^\ell$条轨迹，两个半样本复用同一批轨迹。因此
\begin{align*}
 \E\operatorname{cost}(Z_n^{\db})
 &=n+p_{\rm act} n\sum_{\ell\ge1}\nu_\ell2^\ell
 =n(1+p_{\rm act} \Upsilon_b).
\end{align*}
由$b>1$知$\Upsilon_b<\infty$；对几何级数直接求和便得式~\eqref{eq:level-series-constants}中的闭式。

若$Z_{n,1}^{\db},\ldots,Z_{n,M}^{\db}$独立同分布，则严格无偏与独立性给出
\[
 \E\norm{\widehat g_{M,n}^{\db}-g}_2^2
 =\frac1M\tr\Cov(Z_n^{\db})
 \le\frac{C_{\rm var}(n,p_{\rm act},b)}{M}.
\]
故式~\eqref{eq:estimator-mse}成立。
\end{proof}

\begin{proof}
固定目标与$(n,p_{\rm act},b)$后，定理~\ref{thm:estimator}保证$Z_{n,m}^{\db}$独立同分布，均值为$g$，协方差$\Sigma_n$有限。对任意固定$t\in\R^d$，标量$t^\top(Z_{n,m}^{\db}-g)$独立同分布且方差有限。经典标量中心极限定理给出
\[
 \sqrt M\,t^\top(\widehat g_{M,n}^{\db}-g)
 \rightsquigarrow \mathcal N(0,t^\top\Sigma_nt).
\]
由Cram\'er--Wold装置得到式~\eqref{eq:clt}。

因为$\E\norm{Z_n^{\db}}_2^2<\infty$，矩阵$Z_n^{\db}(Z_n^{\db})^\top$的每个元素都可积。大数定律使$\widehat\Sigma_{M,n}$逐元素依概率收敛于$\Sigma_n$；在固定维度下，这也给出算子范数收敛。若$\Sigma_n$非奇异，则
$\Pp\{\lambda_{\min}(\widehat\Sigma_{M,n})>\lambda_{\min}(\Sigma_n)/2\}\to1$。在该事件上，逆平方根映射于$\Sigma_n$的邻域中有定义且连续。将Slutsky定理用于式~\eqref{eq:clt}，便得到式~\eqref{eq:studentized-clt}。
\end{proof}
